\documentclass{article}
\usepackage[utf8]{inputenc}
\usepackage[margin=0.75in]{geometry}
\usepackage[dvipsnames]{xcolor} % adding colour
\usepackage{amsmath} % adding math equations
\usepackage{natbib} % adding BibTeX support
\usepackage{hyperref}



\begin{document}
	\begin{center}
    
    	% MAKE SURE YOU TAKE OUT THE SQUARE BRACKETS
    
		\LARGE{\textbf{STA380 Project Proposal}} \\
        \vspace{1em}
        \Large{Bootstrap Estimate} \\
        \vspace{1em}
        \normalsize\textbf{Tristan Copplestone, Jacob Gipp, Riclyn Arpilleda, Johnathon Kok} \\
        \normalsize{tristan.copplestone@mail.utoronto.ca, jacob.gipp@mail.utoronto.ca, riclynaarpilleda@mail.utoronto.ca, jonathon.kok@mail.utoronto.ca} \\
        \vspace{1em}
        \normalsize{University of Toronto, Mississauga}
     
	\end{center}
    \begin{normalsize}    
    	\section{Project Topic}
        Computing bootstrap estimates for two descriptive statistics. \\

        \section{Simulation vs. Dataset}
        The results will be from a dataset. We used a dataset that recorded reaction time and aging by using a Nintendo Wii balance board. It was used in the ``Reference data on reaction time and aging using the Nintendo Wii Balance Board: A cross-sectional study of 354 subjects from 20 to 99 years of age" paper.
        
        \section{Project Details}
        We will be computing bootstrap estimates using the dominant/non dominant hand, and dominant/non dominant foot data of the following statistics
        \begin{itemize}
            \item Mean
            \item Interquartile range
        \end{itemize}
        The following outputs will be provided:
        \begin{itemize}
            \item A histogram of the bootstrap distribution of the mean, with a vertical line indicating the true value.
            \item A histogram of the bootstrap distribution of the IQR, with a vertical line indicating the true value.
        \end{itemize}
        
        \section{User Inputs (Shiny Components)}
        Users will be able to modify the following:
        \begin{enumerate}
            \item Setting the random seed.
            \item Choose the number of resamples done for the bootstrap estimates.
            \item Select the information presented: 
                \begin{itemize}
                    \item Gender
                    \item Age range
                    \item Which dataset (dominant/non dominant hand, dominant/non dominant foot)
                \end{itemize}
            \item Be able to choose what information is presented, seeing the mean, or IQR seperately if they wish.
            \item For the plots, be able to modify the colours.
        \end{enumerate}

        \section{Generative AI Statement}
        We did not use generative AI for this checkpoint.

\end{normalsize}

\nocite{*}
\bibliographystyle{abbrvnat}
\bibliography{bibliography}
\vspace{0.8cm}



\end{document}